% !TeX program = xelatex
\documentclass[UTF8,a4paper,zihao=-4,fontset=none]{ctexart}
\setCJKmainfont[BoldFont={LXGW Neo XiHei}]{LXGW Neo ZhiSong}
\setCJKsansfont[AutoFakeBold=2]{LXGW Neo XiHei}
\setCJKmonofont{LXGW Neo XiHei}
\setCJKfamilyfont{zhhei}[AutoFakeBold=2]{LXGW Neo XiHei}
\providecommand{\heiti}{\CJKfamily{zhhei}}
% Shared renderer entrypoint for generated lessonplan.tex files.
% Generated documents should load ctexart, then input this file.

\usepackage{hepingjie_lessonplan}
% Hepingjie lesson-plan overlay coordinates.
% Coordinate origin: top-left corner of an A4 portrait page.
% Units: mm. These coordinates are frozen for templates/lessonplan/hepingjie_blank.pdf.

\newcommand{\LPBackgroundPdf}{templates/lessonplan/hepingjie_blank.pdf}

% Page 1 fixed fields.
\newcommand{\LPDateYearX}{65mm}
\newcommand{\LPDateMonthX}{89mm}
\newcommand{\LPDateDayX}{110mm}
\newcommand{\LPDateY}{38mm}
\newcommand{\LPRecordPageX}{170mm}
\newcommand{\LPRecordPageY}{38mm}

\newcommand{\LPTopicX}{32mm}
\newcommand{\LPTopicY}{49mm}
\newcommand{\LPTopicW}{108mm}
\newcommand{\LPTypeX}{154mm}
\newcommand{\LPTypeY}{49mm}
\newcommand{\LPTypeW}{40mm}

\newcommand{\LPUnitTotalX}{67mm}
\newcommand{\LPUnitTotalY}{61mm}
\newcommand{\LPTopicTotalX}{111mm}
\newcommand{\LPTopicTotalY}{61mm}
\newcommand{\LPCurrentPeriodX}{169mm}
\newcommand{\LPCurrentPeriodY}{61mm}

\newcommand{\LPBodyX}{33mm}
\newcommand{\LPObjectiveY}{74mm}
\newcommand{\LPObjectiveH}{39mm}
\newcommand{\LPKeyPointY}{119mm}
\newcommand{\LPKeyPointH}{15mm}
\newcommand{\LPDifficultyY}{137mm}
\newcommand{\LPDifficultyH}{15mm}
\newcommand{\LPMethodY}{156mm}
\newcommand{\LPMethodH}{15mm}
\newcommand{\LPMediaY}{174mm}
\newcommand{\LPMediaH}{15mm}
\newcommand{\LPBlackboardY}{196mm}
\newcommand{\LPBlackboardH}{73mm}
\newcommand{\LPBodyW}{164mm}

% Teaching process boxes on pages 2-4.
\newcommand{\LPProcessX}{33mm}
\newcommand{\LPProcessY}{31mm}
\newcommand{\LPProcessW}{163mm}
\newcommand{\LPProcessFullH}{249mm}
\newcommand{\LPProcessPageFourH}{160mm}


\usetikzlibrary{calc}
\begin{document}

\LPBackgroundPage{1}{
  \LPCenterBox{\LPDateYearX}{\LPDateY}{12mm}{5mm}{}
  \LPCenterBox{\LPDateMonthX}{\LPDateY}{10mm}{5mm}{}
  \LPCenterBox{\LPDateDayX}{\LPDateY}{10mm}{5mm}{}
  \LPCenterBox{\LPRecordPageX}{\LPRecordPageY}{10mm}{5mm}{1}
  \LPHeaderCell{\LPTopicX}{\LPTopicY}{\LPTopicW}{8mm}{几何图形初步章末小结与复习}
  \LPHeaderCell{\LPTypeX}{\LPTypeY}{\LPTypeW}{8mm}{章末复习课}
  \LPCenterBox{\LPUnitTotalX}{\LPUnitTotalY}{12mm}{5mm}{11}
  \LPCenterBox{\LPTopicTotalX}{\LPTopicTotalY}{12mm}{5mm}{6}
  \LPCenterBox{\LPCurrentPeriodX}{\LPCurrentPeriodY}{12mm}{5mm}{6}
  \LPTextBox{\LPBodyX}{\LPObjectiveY}{\LPBodyW}{\LPObjectiveH}{\LPPlainLine{说出并解释复习几何图形、直线射线线段和角，使用本节规范术语表达。}\LPPlainLine{按规范步骤完成几何图形知识结构，并说明关键依据。}\LPPlainLine{判断并订正与“图形语言、符号语言和文字语言的转换”有关的常见错误。}\LPPlainLine{在变式或实际情境中运用本节方法，完整写出过程和结论。}}
  \LPTextBox{\LPBodyX}{\LPKeyPointY}{\LPBodyW}{\LPKeyPointH}{几何图形知识结构；线段和角的综合}
  \LPTextBox{\LPBodyX}{\LPDifficultyY}{\LPBodyW}{\LPDifficultyH}{图形语言、符号语言和文字语言的转换}
  \LPTextBox{\LPBodyX}{\LPMethodY}{\LPBodyW}{\LPMethodH}{结构梳理；题组训练；错因诊断}
  \LPTextBox{\LPBodyX}{\LPMediaY}{\LPBodyW}{\LPMediaH}{PPT；黑板；反馈卡；步骤卡}
  \LPTextBox{\LPBodyX}{\LPBlackboardY}{\LPBodyW}{\LPBlackboardH}{\LPBoardTitle{几何图形初步章末小结与复习}\LPBoardLine{1}{核心：先读点线角名称，标已知条件，确定点或射线顺序，再列线段或角的和差关系。}\LPBoardLine{2}{方法：先标条件和顺序，再用定义、基本事实与和差关系列式。}\LPBoardLine{3}{例题：$M$ 为 $AB$ 中点，$AB=14$ cm，$C$ 在线段 $MB$ 上且 $MC=2$ cm，求 $AC$。 答：$AM=7$ cm，$AC=AM+MC=9$ cm。}\LPBoardLine{4}{易错：错解：仅凭示意图断定 $M$ 是中点。 纠正：错误；必须有文字条件或相等线段标记。}\LPBoardLine{5}{检查：条件、依据、步骤、符号或单位逐项核对。}}
}

\LPBackgroundPage{2}{
  \LPProcessBox{\LPProcessX}{\LPProcessY}{\LPProcessW}{\LPProcessFullH}{\LPProcessRow{知识诊断}{\LPTeacherPart{问题}{本章立体与平面图形、点线面体、线段和角怎样组成知识结构？}\LPTeacherPart{组织}{出示题目或情境，先让学生独立判断，再请两名学生说明依据。}\LPTeacherPart{说明}{由基本图形抽象点线面体，再研究直线射线线段，最后研究角的表示、运算和关系。}\LPTeacherPart{纠错}{若学生只报结论，追问基准、条件或运算依据，要求补成完整句。}\LPTeacherPart{归纳}{把情境中的信息转化为数学对象和关系。}\LPTeacherPart{意图}{用具体问题激活先修经验，并明确本节需要解决的核心矛盾。}}{\LPStudentPart{活动}{独立观察或计算，圈出关键信息后口头说明。}\LPStudentPart{预期}{由基本图形抽象点线面体，再研究直线射线线段，最后研究角的表示、运算和关系。}}{5}
\LPProcessRow{结构梳理}{\LPTeacherPart{问题}{几何综合题如何读图和表达？}\LPTeacherPart{组织}{组织观察、比较或试算，板书学生给出的共同点，并逐项核对术语。}\LPTeacherPart{说明}{先读点线角名称，标已知条件，确定点或射线顺序，再列线段或角的和差关系。}\LPFigureBlock{
\begin{tikzpicture}[scale=.68,baseline=(current bounding box.center)]
  \draw (0,0)--(5,0); \foreach \x/\t in {0/A,2.5/M,4/C,5/B}{\fill(\x,0)circle(1.2pt)node[above]{$\t$};}
  \node[below]at(2.5,-.25){\scriptsize 线段和差与中点};
  \begin{scope}[xshift=7cm]
    \draw[->](0,0)--(3,0)node[right]{$A$};\draw[->](0,0)--(2.2,1.1)node[right]{$B$};\draw[->](0,0)--(.7,2.5)node[above]{$C$};
    \node[left]at(0,0){$O$};\draw (.5,0)arc(0:27:.5);\draw (27:.5)arc(27:74:.5);
  \end{scope}
\end{tikzpicture}
}\LPTeacherPart{例题}{基础例题：$M$ 为 $AB$ 中点，$AB=14$ cm，$C$ 在线段 $MB$ 上且 $MC=2$ cm，求 $AC$。 规范解答：$AM=7$ cm，$AC=AM+MC=9$ cm。}\LPTeacherPart{对应练习}{即时练习：说出直线、射线、线段端点数。 答案：0、1、2。}\LPTeacherPart{纠错}{用反例检查是否真正满足条件，提醒按“先标条件和顺序，再用定义、基本事实与和差关系列式”说明。}\LPTeacherPart{归纳}{先标条件和顺序，再用定义、基本事实与和差关系列式}\LPTeacherPart{意图}{让核心结论由可检验的观察或运算形成，而不是直接记忆。}}{\LPStudentPart{活动}{在图形、算式或表格中标注条件，与同桌归纳共同特征。}\LPStudentPart{预期}{先读点线角名称，标已知条件，确定点或射线顺序，再列线段或角的和差关系。}}{7}
\LPProcessRow{典型题组}{\LPTeacherPart{问题}{解答“$OB$ 平分 $\angle AOC=80^\circ$，$OD$ 在 $\angle BOC$ 内且 $\angle BOD=15^\circ$，求 $\angle DOC$。”前应先确定什么？}\LPTeacherPart{组织}{教师按题意、依据、步骤和结论顺序板演；学生在每一步旁标注作用。}\LPTeacherPart{说明}{规范解答：$\angle BOC=40^\circ$，所以 $\angle DOC=25^\circ$。}\LPTeacherPart{例题}{变式例题：$OB$ 平分 $\angle AOC=80^\circ$，$OD$ 在 $\angle BOC$ 内且 $\angle BOD=15^\circ$，求 $\angle DOC$。 解答：$\angle BOC=40^\circ$，所以 $\angle DOC=25^\circ$。}\LPTeacherPart{对应练习}{对应练习：求 $35^\circ$ 的补角。 答案：$145^\circ$。}\LPTeacherPart{纠错}{若一步跳到结论，要求补写关键关系或中间步骤，并用原题条件检验。}\LPTeacherPart{归纳}{解题流程：先标条件和顺序，再用定义、基本事实与和差关系列式。}\LPTeacherPart{意图}{用完整例题建立可模仿的书写顺序和检查方法。}}{\LPStudentPart{活动}{先独立完成，再与板演对照步骤、符号、单位或图形标记。}\LPStudentPart{预期}{能得到 $\angle BOC=40^\circ$，所以 $\angle DOC=25^\circ$。，并说出关键依据。}}{8}}
}

\LPBackgroundPage{3}{
  \LPProcessBox{\LPProcessX}{\LPProcessY}{\LPProcessW}{\LPProcessFullH}{\LPProcessRow{方法迁移}{\LPTeacherPart{问题}{点 $A,B$ 之间有直线路径 $AB=6$ km 和折线路径 $AC+CB=8$ km，选择较短路线并说明依据。}\LPTeacherPart{组织}{要求学生先写条件关系或示意，再完成计算并解释结果；抽取两种方法比较。}\LPTeacherPart{说明}{选择线段 $AB$；两点之间线段最短。}\LPTeacherPart{例题}{应用例题：点 $A,B$ 之间有直线路径 $AB=6$ km 和折线路径 $AC+CB=8$ km，选择较短路线并说明依据。 解答：选择线段 $AB$；两点之间线段最短。}\LPTeacherPart{对应练习}{即时变式：若 $AB=10$ cm、$M$ 为中点，求 $AM$。 答案：5 cm。}\LPTeacherPart{纠错}{若答案脱离条件或单位，要求回到题干逐项对应并作合理性检查。}\LPTeacherPart{归纳}{先识别结构，再选择方法，最后解释结果。}\LPTeacherPart{意图}{把核心方法迁移到不同表示方式或实际情境中。}}{\LPStudentPart{活动}{独立列式或作图，小组内互查条件、步骤和答语。}\LPStudentPart{预期}{选择线段 $AB$；两点之间线段最短。}}{6}
\LPProcessRow{易错诊断}{\LPTeacherPart{问题}{错解：仅凭示意图断定 $M$ 是中点。}\LPTeacherPart{组织}{展示错解，要求学生只圈第一处错误，写出依据后再完整订正。}\LPTeacherPart{说明}{错误；必须有文字条件或相等线段标记。}\LPTeacherPart{例题}{易错辨析：错解：仅凭示意图断定 $M$ 是中点。}\LPTeacherPart{对应练习}{纠错要求：写出第一处错误，并按“先标条件和顺序，再用定义、基本事实与和差关系列式”重做。参考：错误；必须有文字条件或相等线段标记。}\LPTeacherPart{纠错}{不接受只写“错”或只改答案；必须指出错误对象、正确规则和订正过程。}\LPTeacherPart{归纳}{纠错顺序：定位错误、说明依据、规范订正、回题检验。}\LPTeacherPart{意图}{利用典型错误暴露概念边界、符号习惯或条件遗漏。}}{\LPStudentPart{活动}{独立找错、写依据、完成订正，再与同桌互评理由是否完整。}\LPStudentPart{预期}{错误；必须有文字条件或相等线段标记。}}{5}
\LPProcessRow{综合训练}{\LPTeacherPart{问题}{三道练习分别考查哪个要点，先做哪一道能最快暴露问题？}\LPTeacherPart{组织}{组织限时题组：①说出直线、射线、线段端点数。 ②求 $35^\circ$ 的补角。 ③若 $AB=10$ cm、$M$ 为中点，求 $AM$。；完成后按答案自查。}\LPTeacherPart{说明}{答案：①0、1、2。 ②$145^\circ$。 ③5 cm。}\LPTeacherPart{对应练习}{机动题：$M$ 为 $AB$ 中点，$AB=14$ cm，$C$ 在线段 $MB$ 上且 $MC=2$ cm，求 $AC$。 学有余力者换一种方法说明；答案要点：$AM=7$ cm，$AC=AM+MC=9$ cm。}\LPTeacherPart{纠错}{正确率低于 70\% 时暂停机动题，按核心方法示范一道，再让学生订正同类错题。}\LPTeacherPart{归纳}{先标条件和顺序，再用定义、基本事实与和差关系列式}\LPTeacherPart{意图}{集中检查方法稳定性，并为反馈测试前的即时补救提供依据。}}{\LPStudentPart{活动}{限时完成三题，按概念、步骤或应用给错因分类，并订正一题。}\LPStudentPart{预期}{能够依据“先标条件和顺序，再用定义、基本事实与和差关系列式”完成题组并说清错因。}}{4}}
}

\LPBackgroundPage{4}{
  \LPProcessBox{\LPProcessX}{\LPProcessY}{\LPProcessW}{\LPProcessPageFourH}{\LPProcessRow{当堂检测}{\LPTeacherPart{问题}{四道反馈题中，哪一题最能检验本节核心方法？}\LPTeacherPart{组织}{投放 10 分反馈卡：①2 分，几何综合题如何读图和表达？ ②2 分，求 $35^\circ$ 的补角。 ③2 分，错解：仅凭示意图断定 $M$ 是中点。 ④4 分，点 $A,B$ 之间有直线路径 $AB=6$ km 和折线路径 $AC+CB=8$ km，选择较短路线并说明依据。}\LPTeacherPart{说明}{参考答案：①先读点线角名称，标已知条件，确定点或射线顺序，再列线段或角的和差关系。 ②$145^\circ$。 ③错误；必须有文字条件或相等线段标记。 ④选择线段 $AB$；两点之间线段最短。}\LPTeacherPart{纠错}{公布答案后学生自评并举牌统计：9 分及以上完成提高任务；7—8 分订正；7 分以下接受针对性补练。}\LPTeacherPart{归纳}{达标标准：10 分制 9 分及以上完全达标，7—8 分基本达标，7 分以下补练。}\LPTeacherPart{意图}{用概念、技能、纠错和综合四类任务覆盖本节目标。}}{\LPStudentPart{活动}{独立完成，按答案自评；把错误标为概念、步骤、条件或表达问题。}\LPStudentPart{预期}{能完成四类题并用核心术语说明至少一道题的依据。}}{6}
\LPProcessRow{分层作业}{\LPTeacherPart{问题}{本节研究了什么、关键步骤是什么、最容易错在哪里？}\LPTeacherPart{组织}{请学生各用一句话回答三个问题，教师据此补全板书并布置分层作业。}\LPTeacherPart{说明}{A 组：①说出直线、射线、线段端点数。 ②求 $35^\circ$ 的补角。 ③若 $AB=10$ cm、$M$ 为中点，求 $AM$。 ④错解：仅凭示意图断定 $M$ 是中点。；答案：①0、1、2。 ②$145^\circ$。 ③5 cm。 ④错误；必须有文字条件或相等线段标记。。B 组：①$OB$ 平分 $\angle AOC=80^\circ$，$OD$ 在 $\angle BOC$ 内且 $\angle BOD=15^\circ$，求 $\angle DOC$。 ②点 $A,B$ 之间有直线路径 $AB=6$ km 和折线路径 $AC+CB=8$ km，选择较短路线并说明依据。；答案要点：①$\angle BOC=40^\circ$，所以 $\angle DOC=25^\circ$。 ②选择线段 $AB$；两点之间线段最短。。C 组选做：围绕“几何图形初步章末小结与复习”自编一道含易错条件的题并给出解答与检查方法。}\LPTeacherPart{纠错}{作业答案必须包含必要步骤或依据；只写结果的题次日先口述方法再补写。}\LPTeacherPart{归纳}{核心方法：先标条件和顺序，再用定义、基本事实与和差关系列式。课后反思区域由任课教师课后填写。}\LPTeacherPart{意图}{由学生完成结构化回顾，把课堂方法迁移到分层课后任务。}}{\LPStudentPart{活动}{说出一个结论、一个步骤和一个易错点，记录 A、B、C 三组作业。}\LPStudentPart{预期}{能用“先标条件和顺序，再用定义、基本事实与和差关系列式”概括本节方法，并知道如何自查。}}{4}}
}
\end{document}
